% ============================================================
%  Artigo 01 — A Comprehensive Benchmark on Classical PID Tuning Strategies
%  IC/PROIC-UESC — Template para Overleaf
%  Baseado no estilo IEEE Transactions
% ============================================================

\documentclass[journal]{IEEEtran}

% ---- Pacotes essenciais ----
\usepackage[T1]{fontenc}
\usepackage[utf8]{inputenc}
\usepackage[english]{babel}
\usepackage{amsmath, amssymb, amsfonts}
\usepackage{graphicx}
\usepackage{booktabs}
\usepackage{hyperref}
\usepackage{siunitx}
\usepackage{algorithm}
\usepackage{algpseudocode}
\usepackage{cite}

% ---- Metadados ----
\title{A Comprehensive Benchmark on Classical PID Tuning Strategies for Highly Non-linear CSTR: Robustness, Actuator Saturation, and Multi-Objective Trade-offs}

\author{%
  Gabriel Marques~\IEEEmembership{Student Member, UESC},
  Orientador~\IEEEmembership{Professor, UESC}%
  \thanks{Apoio: CNPq / FAPESB / PROIC-UESC. Edital 34/2026.}%
}

\begin{document}

\maketitle

% ============================================================
\begin{abstract}
This paper presents an exhaustive comparative benchmark of classical Proportional-Integral-Derivative (PID) tuning methodologies applied to a simulated non-isothermal Continuous Stirred Tank Reactor (CSTR) characterized by severe Arrhenius non-linearities and realistic actuator saturation dynamics. Four classical tuning paradigms—Ziegler--Nichols (ZN), Cohen--Coon (CC), Skogestad Simple Internal Model Control (SIMC), and Constrained Numerical Optimization (ITAE with maximum sensitivity $M_s \le 1.6$)—are systematically evaluated across nominal setpoint tracking, unmeasured load disturbance rejection (feed temperature surge $\Delta T_f = +5$~K and concentration surge $\Delta C_{Af} = +20\%$), heat exchanger thermal fouling ($UA$ degradation from $100\%$ down to $60\%$), sensor measurement noise, and multi-operating point transitions. The results reveal that aggressive empirical formulas (ZN and CC) induce persistent valve saturation, thermal runaway risk, and destructive limit cycles. In contrast, Skogestad SIMC and Constrained Optimal ITAE maintain superior damping ($M_p \le 12\%$, $IAE < 1.8$) and robustly reject severe load disturbances while safeguarding the actuator. Furthermore, a Pareto frontier analysis maps the fundamental trade-off between tracking accuracy ($IAE$) and control effort ($TV$), establishing a rigorous empirical baseline for forthcoming investigations into adaptive neural and agent-based supervisory control.
\end{abstract}

\begin{IEEEkeywords}
PID control, CSTR, chemical process control, Skogestad SIMC, actuator saturation, anti-windup, thermal fouling, Pareto optimization, robustness benchmark
\end{IEEEkeywords}

% ============================================================
\section{Introduction}
\label{sec:intro}
% ============================================================
%  Seção 1: Introduction — Rascunho
%  Preencher / refinar conforme revisão bibliográfica avança
% ============================================================

% TODO: Revisar e expandir
Proportional-Integral-Derivative (PID) controllers remain the dominant control 
strategy in the chemical process industry, accounting for over 90\% of regulatory 
control loops \cite{astrom2006pid}. Despite their widespread adoption, classical 
fixed-gain PID controllers face significant challenges in highly nonlinear processes 
such as Continuous Stirred Tank Reactors (CSTRs), where Arrhenius kinetics introduce 
exponential temperature-reaction coupling and the risk of thermal runaway 
\cite{seborg2010process}.

Recent advances in machine learning, particularly deep neural networks, have opened 
new avenues for adaptive control that can dynamically adjust PID gains based on 
real-time process states \cite{pid_neural_survey}. Furthermore, the emergence of 
large language model-based agent frameworks such as LangGraph enables the orchestration 
of specialized autonomous agents into coherent control pipelines, offering modular, 
interpretable, and extensible control architectures.

This paper makes the following contributions:
\begin{enumerate}
    \item A systematic comparison of three classical PID tuning methods (Ziegler–Nichols, 
    Cohen–Coon, and IMC) on a simulated CSTR;
    \item A neural network-based adaptive PID controller (Neuro-PID) implemented 
    in PyTorch that dynamically adjusts gains based on process state;
    \item A multi-agent orchestration framework using LangGraph that integrates 
    telemetry, surrogate modeling, neural tuning, and safety verification;
    \item Quantitative performance evaluation using IAE, ITAE, and TV metrics 
    under standardized test scenarios.
\end{enumerate}


% ============================================================
\section{Mathematical Modeling of the CSTR}
\label{sec:model}
% ============================================================
%  Section 2: Mathematical Modeling of the Non-linear CSTR and Controllability Limits
% ============================================================

\subsection{Non-linear Physical Plant Model}
A continuous stirred tank reactor (CSTR) conducting an irreversible, liquid-phase exothermic reaction $A \xrightarrow{k} B$ is considered \cite{luyben2007chemical, seborg2010process}. Rather than adopting the frequent simplification of instantaneous jacket dynamics, the cooling jacket is modeled explicitly as a dynamic lumped thermal system coupled through non-linear heat convection across the reactor wall \cite{uppal1974dynamic}.

The dynamic behavior of the process is governed by three coupled non-linear ordinary differential equations (ODEs) derived from dynamic species mass and thermal energy balances:
\begin{align}
    \frac{dC_A}{dt} &= \frac{q}{V}(C_{Af} - C_A) - r_A(C_A, T) \label{eq:cstr_ca} \\
    \frac{dT}{dt}   &= \frac{q}{V}(T_f - T) + \frac{(-\Delta H)}{\rho C_p} r_A(C_A, T) - \frac{UA}{V \rho C_p}(T - T_j) \label{eq:cstr_t} \\
    \frac{dT_j}{dt} &= \frac{q_j}{V_j}(T_{jf} - T_j) + \frac{UA}{V_j \rho_j C_{pj}}(T - T_j) \label{eq:cstr_tj}
\end{align}
where $C_A$ denotes the concentration of reactant $A$ in the reactor, $T$ is the bulk reactor temperature, and $T_j$ is the cooling jacket temperature. The non-linear reaction rate $r_A(C_A, T)$ obeys Arrhenius kinetics:
\begin{equation}
    r_A(C_A, T) = k_0 \exp\left(-\frac{E}{R\,T}\right) C_A
    \label{eq:arrhenius}
\end{equation}
The physical state vector is $\mathbf{x} = [C_A, T, T_j]^T \in \mathbb{R}^3$, the manipulated input is the coolant volumetric flow rate $u(t) = q_j(t) \in \mathbb{R}$, and the measured process output is the reactor core temperature $y(t) = T(t)$.

\subsection{Actuator Physical Constraints and Slew-Rate Dynamics}
In physical chemical plants, control valves are constrained by finite physical stroke and actuator motor torque limits \cite{hippe2006windup, tarbouriech2009antiwindup}. The manipulated coolant flow $q_j(t)$ is bounded by hard amplitude saturation and valve travel velocity (slew-rate) limits:
\begin{align}
    u_{\min} &\le q_j(t) \le u_{\max} \label{eq:actuator_amplitude} \\
    \left|\frac{dq_j}{dt}\right| &\le S_{\max} \label{eq:actuator_slew}
\end{align}
where $u_{\min} = 0.0$~L/min represents a fully closed valve, $u_{\max} = 300.0$~L/min corresponds to full valve opening, and $S_{\max} = 80.0$~L/min$^2$ models the electro-pneumatic actuator positioning speed. Slew-rate clamping introduces an effective dynamic dead time when sudden cooling surges are commanded, severely exacerbating the risk of overshoot.

\subsection{Nominal Operating Parameters and Equilibrium}
The nominal physico-chemical parameters are compiled in Table~\ref{tab:cstr_parameters}. Under the nominal steady-state coolant flow rate $q_{j,ss} = 100.0$~L/min, the non-linear algebraic system $\mathbf{f}(\mathbf{x}_{ss}, q_{j,ss}) = \mathbf{0}$ yields the nominal operating point:
\begin{equation}
    \mathbf{x}_{ss} = \begin{bmatrix} C_{A,ss} \\ T_{ss} \\ T_{j,ss} \end{bmatrix} = \begin{bmatrix} 0.0502\text{ mol/L} \\ 396.65\text{ K} \\ 348.33\text{ K} \end{bmatrix}
    \label{eq:nominal_ss}
\end{equation}
This equilibrium corresponds to a high reactant conversion of $X_A = (C_{Af} - C_{A,ss})/C_{Af} \approx 95\%$. At this elevated conversion, the reactor operates near maximum heat release, where small positive temperature perturbations trigger exponential acceleration of the reaction rate.

\begin{table}[htbp]
\centering
\caption{Nominal physical and kinetic parameters of the simulated CSTR.}
\label{tab:cstr_parameters}
\begin{tabular}{llrl}
\toprule
Parameter & Symbol & Value & Unit \\
\midrule
Reactor volume & $V$ & $100.0$ & L \\
Cooling jacket volume & $V_j$ & $20.0$ & L \\
Process volumetric flow rate & $q$ & $100.0$ & L/min \\
Feed concentration & $C_{Af}$ & $1.0$ & mol/L \\
Feed temperature & $T_f$ & $350.0$ & K \\
Coolant feed temperature & $T_{jf}$ & $300.0$ & K \\
Pre-exponential factor & $k_0$ & $7.2 \times 10^{10}$ & min$^{-1}$ \\
Activation energy parameter & $E/R$ & $8750.0$ & K \\
Heat of reaction & $-\Delta H$ & $5.0 \times 10^4$ & cal/mol \\
Fluid heat capacity product & $\rho C_p$ & $500.0$ & cal/(L$\cdot$K) \\
Coolant heat capacity product & $\rho_j C_{pj}$ & $500.0$ & cal/(L$\cdot$K) \\
Overall heat transfer coefficient & $UA$ & $5.0 \times 10^4$ & cal/(min$\cdot$K) \\
Nominal coolant flow & $q_{j,ss}$ & $100.0$ & L/min \\
\bottomrule
\end{tabular}
\end{table}

\subsection{Thermal Controllability Boundaries and Van Heerden Analysis}
\label{subsec:thermal_controllability}
The controllability of exothermic CSTRs is fundamentally bounded by the balance between non-linear heat generation $Q_g(T)$ and dynamic heat removal $Q_r(T, q_j)$ \cite{aris1958analysis, uppal1974dynamic}:
\begin{align}
    Q_g(T) &= (-\Delta H) V r_A(C_A(T), T) \label{eq:q_gen} \\
    Q_r(T, q_j) &= q \rho C_p (T - T_f) + UA (T - T_j(q_j)) \label{eq:q_rem}
\end{align}
According to the classic Van Heerden stability criterion \cite{seborg2010process, bequette2003process}, open-loop thermal runaway occurs whenever the temperature sensitivity of heat generation exceeds that of heat removal:
\begin{equation}
    \left. \frac{\partial Q_g}{\partial T} \right|_{\mathbf{x}_{ss}} > \left. \frac{\partial Q_r}{\partial T} \right|_{\mathbf{x}_{ss}}
    \label{eq:van_heerden}
\end{equation}
At nominal steady state, the system possesses an open-loop unstable eigenvalue ($\lambda_1 = +0.284$~min$^{-1}$), requiring active feedback stabilization.

Furthermore, physical actuator saturation imposes an absolute limit on heat removal. Even at full valve capacity ($q_j = u_{\max} = 300$~L/min), the maximum achievable cooling rate is strictly bounded by:
\begin{equation}
    Q_{r, \max}(T) = q \rho C_p (T - T_f) + UA (T - T_{jf})
    \label{eq:max_cooling_limit}
\end{equation}
When the overall heat transfer capability degrades (e.g., due to polymer fouling on the heat exchange surface such that $UA < UA_{\text{crit}} \approx 0.65 UA_0$), $Q_g(T) > Q_{r, \max}(T)$ over a finite temperature interval. Under such conditions, the plant enters an unrecoverable thermal runaway regime regardless of the control law employed. This analytical boundary establishes the physical foundation for the fouling and load disturbance stress tests evaluated herein.


% ============================================================
\section{Classical PID Tuning Architectures}
\label{sec:pid}
% Placeholder sections — completar conforme desenvolvimento das Fases 3-6
% TODO: Preencher com conteúdo técnico após execução das fases



% ============================================================
\section{Simulation Results and Robustness Benchmark}
\label{sec:results}
% TODO — Preencher com resultados da Fase 6


% ============================================================
\section{Conclusion and Future Work}
\label{sec:conclusion}
% ============================================================
%  Section 7: Conclusion and Future Directions
% ============================================================

This paper presented an exhaustive, rigorous computational benchmark evaluating the physical controllability boundaries and robustness of classical PID tuning strategies on a highly non-linear, non-isothermal Continuous Stirred Tank Reactor (CSTR) subject to realistic actuator amplitude saturation and slew-rate dynamics. Four tuning paradigms—Ziegler--Nichols (ZN), Cohen--Coon (CC), Skogestad Simple Internal Model Control (SIMC), and Constrained Numerical Optimization (ITAE with an explicit $\mathcal{H}_\infty$ maximum sensitivity constraint $M_s \le 1.6$)—were systematically investigated across nominal setpoint tracking, unmeasured exothermic load disturbances, heat exchanger surface fouling, sensor noise, and multi-operating point transitions.

The investigation yields several fundamental conclusions for process control theory and industrial practice:
\begin{enumerate}
    \item \textbf{Predictable Collapse of Empirical Reaction-Curve Rules:} Classical ZN and CC tuning formulas prescribe excessively aggressive proportional gains ($K_p \approx -50$) when applied to lag-dominant systems ($\theta/\tau \approx 0.085$). Under realistic actuator bounds, this unconstrained aggressiveness immediately drives the cooling valve into persistent saturation, inducing violent limit cycling and catastrophic temperature excursions ($\Delta T > 43$~K) during exothermic disturbances. This confirms that empirical formulas formulated for unconstrained linear plants cannot be safely deployed on lag-dominant exothermic reactors.
    
    \item \textbf{Superiority of $\mathcal{H}_\infty$ Robust Design and Model-Based Control:} Skogestad SIMC ($\tau_c = 0.40$~min) and Constrained Optimal ITAE ($M_s \le 1.6$) demonstrate exceptional closed-loop resilience, achieving well-damped tracking ($\text{IAE} < 1.8$, $M_p \le 12\%$) without actuator saturation, and attenuating sequential unmeasured thermal and concentration shocks with peak deviations confined to under $7.4$~K.
    
    \item \textbf{Analytical and Physical Controllability Limits:} Parametric fouling tests revealed that both robust strategies preserve stability down to $70\% UA$. At $60\% UA$, the maximum cooling capacity of the jacket is physically incapable of rejecting the exothermic reaction heat, delineating the absolute thermodynamic controllability boundary where linear feedback control authority collapses.
    
    \item \textbf{Authentic Pareto Frontier and Sensitivity Boundaries:} Continuous sweeps of the closed-loop design parameter $\tau_c$ and scalarized multi-objective optimization rigorously mapped the non-linear Pareto trade-off between tracking accuracy ($\text{IAE}$) and actuator physical wear ($\text{TV}$). The analysis demonstrates that selecting $\tau_c < 0.10$~min yields diminishing returns in tracking performance while driving peak sensitivity beyond $M_s > 2.0$ and drastically escalating valve chattering. Enforcing $M_s \le 1.6$ provides a reliable design guideline that guarantees robust stability margins while preventing premature actuator wear.
    
    \item \textbf{Foundation for Adaptive Neural and Supervisory Architectures:} Wide-range multistep tracking across disparate reaction regimes confirmed that fixed-gain linear controllers inevitably suffer performance degradation when operating away from the linearization design point. The quantitative metrics, controllability boundaries, and Pareto limits established herein provide an authoritative benchmark against which forthcoming adaptive neural PID and autonomous multi-agent supervisory architectures will be evaluated.
\end{enumerate}


% ============================================================
\bibliographystyle{IEEEtran}
\bibliography{references}

\end{document}
